% Figures du lot « espaces latents » du Tome III : III-7.6 (l'espace latent des états de la couche 23),
% III-9.1 (d) (une rangée de 32 poids avant et après la mutation synonyme), III-9.7 (l'erreur de sortie d'un
% expert dans les directions de ses entrées). Gardé par \ifdefined\GrapheLatentsEtats.
% Données : \Racine/figures/tome3/donnees/latents-*.csv et latents-chiffres.tex (chaque nombre écrit dans
% une légende ou une annotation), écrits par memoire/outils/tome3_latents.py depuis mesures/genome/tome3/
% latents/ (outils/genome/tome3/latents_*.py). Styles et couleurs propres au lot préfixés « lat ».
%
% Second flux (2026-10-27) : palette et convention de légende du Tome III (palette-tome3.tex,
% legende-tome3.tex ; memoire/tome3/CONVENTIONS-figures.md) ; graves du juge de design G1 (couleurs du
% dictionnaire), G2 (III-7.6 en petits multiples) et D20 (III-9.7 lue par la fonction d'appui). Détail et
% décisions : memoire/tome3/REPRISE-figures2-latents.md.
%
% Pour chaque figure <X> (LatentsEtats, LatentsRangee, LatentsNuage) : \Graphe<X> (le dessin), et, par la
% convention G9, \Titre<X>, \Legende<X>, \Source<X>. Emploi au chapitre :
%   \begin{figure}\centering\GrapheLatentsEtats\LegendeFigure{LatentsEtats}\label{...}\end{figure}

\ifdefined\GrapheLatentsEtats\else
\ifdefined\LegendeTomeTrois\else\input{\Racine/figures/tome3/legende-tome3}\fi
\newcommand{\LatDonnees}{\Racine/figures/tome3/donnees}
% Un pourcentage après une macro de nombre : « 22.0% » en anglais, « 22,0 % » en français.
\newcommand{\LatPc}{\EF{\%}{\,\%}}
\input{\LatDonnees/latents-chiffres.tex}

% --- les couleurs : celles du dictionnaire du Tome III, sous leur nom ---------------------------------
% III-7.6 : une seule teinte pour l'expert d'un panneau (la place du panneau dit l'expert, la couleur ne
% porte plus rien seule) ; les états de petite marge en encre ; le reste du nuage en gris.
\colorlet{latExpert}{CouleurCategorieA}
\colorlet{latFragile}{Encre}
\colorlet{latFond}{CouleurAutres}
\colorlet{latCarte}{black!30}
% III-9.1 (d) et III-9.7 : la frontière et le synonyme du dictionnaire ; les entrées en gris moyen.
\colorlet{latEntree}{black!35}

\pgfplotsset{
  % une carte d'états : axes égaux, sans grille
  lat/carte/.style={t3 graphe, scale only axis, axis equal image, grid=none,
    axis line style={black!30}, tick style={black!30}, tick label style={font=\scriptsize},
    label style={font=\scriptsize}},
  lat/muet/.style={xtick=\empty, ytick=\empty},
  % le nuage d'une carte : gris, et les états de petite marge en encre par-dessus
  lat/nuage gris/.style={scatter, only marks, mark=*, scatter src=explicit symbolic,
    scatter/classes={autre={mark size=0.42pt, draw=latCarte, fill=latCarte, line width=0.1pt},
                     petite={mark size=0.62pt, draw=latFragile, fill=latFragile, line width=0.1pt}}},
  % troisième passe (DN13 du juge de design) : dans (c), un état fragile est un anneau d'encre autour d'un point de
  % l'expert, plus un point noir plein ; en gris et à petite échelle, le noir plein se fondait dans le bleu
  lat/nuage expert/.style={scatter, only marks, mark=*, scatter src=explicit symbolic,
    scatter/classes={autre={mark size=0.55pt, draw=latExpert, fill=latExpert, line width=0.1pt},
                     petite={mark size=0.95pt, draw=latFragile, fill=latExpert!45, line width=0.35pt}}},
  lat/fond/.style={only marks, mark=*, mark size=0.32pt, mark options={draw=latFond, fill=latFond, line width=0.1pt}},
  % le nom des experts, posé sur la carte
  lat/noms/.style={only marks, mark=none, point meta=explicit symbolic, nodes near coords,
    nodes near coords style={t3 etiquette, font=\sffamily\scriptsize\bfseries, anchor=center,
      fill opacity=0.8}},
}
% Une marque se trace toujours avec son contour (draw=none est ignoré pour les marques sous le style
% memoire : le contour sort en noir) : contour et fond de la même couleur, contour de 0,1 pt.
% Une pastille de clé : #1 couleur, #2 rayon (pt).
\newcommand{\LatPastille}[2]{\tikz[baseline=-0.55ex]\fill[#1] (0,0) circle (#2pt);}

% =====================================================================================================
% III-7.6. L'espace latent des 6 132 états de la couche 23.
% Un panneau des petits multiples. #1 : numéro du groupe (1 à 8) ; #2 : son nom (Un à Huit) ;
% #3 : colonne (0 à 3) ; #4 : ligne (0 ou 1).
\newcommand{\LatMultiple}[4]{%
\begin{axis}[lat/carte, lat/muet, width=2.55cm, height=2.55cm, at={(#3*2.95cm, -5.8cm-#4*2.75cm)},
  anchor=north west, axis line style={black!20}]
\addplot[lat/fond] table[x=x, y=y] {\LatDonnees/latents-etats-tsne.csv};
\addplot[lat/nuage expert] table[x=x, y=y, meta=marge] {\LatDonnees/latents-etats-tsne-g#1.csv};
% l'expert et sa part, dans le coin haut gauche du panneau
\node[t3 etiquette, anchor=north west, inner sep=1.2pt] at (rel axis cs:0.02,0.98)
  {\textbf{\csname LatGroupe#2\endcsname}\enspace{\color{black!55}\csname LatGroupe#2Part\endcsname\LatPc}};
\end{axis}}
\newcommand{\GrapheLatentsEtats}{%
\begin{tikzpicture}
% (a) l'ACP
\begin{axis}[lat/carte, name=latacp, width=4.3cm, height=4.3cm, at={(0,0)}, anchor=north west,
  title={\EF{(a) PCA of the states}{(a) ACP des états}},
  xlabel={\EF{PC1 (\LatAcpVarUn\% of the variance)}{CP1 (\LatAcpVarUn\,\% de la variance)}},
  ylabel={\EF{PC2 (\LatAcpVarDeux\%)}{CP2 (\LatAcpVarDeux\,\%)}}, ylabel shift=-3pt]
\addplot[lat/nuage gris] table[x=x, y=y, meta=marge] {\LatDonnees/latents-etats-acp.csv};
\addplot[t3 rappel, quiver={u=\thisrow{acpx}-\thisrow{acplx}, v=\thisrow{acpy}-\thisrow{acply}}, forget plot]
  table[x=acplx, y=acply] {\LatDonnees/latents-etats-reperes.csv};
\addplot[lat/noms] table[x=acplx, y=acply, meta=expert] {\LatDonnees/latents-etats-reperes.csv};
\end{axis}
% (b) la t-SNE
\IfFileExists{\LatDonnees/latents-etats-tsne.csv}{%
\begin{axis}[lat/carte, lat/muet, name=lattsne, width=4.3cm, height=4.3cm, at={(7.15cm,0)}, anchor=north west,
  title={\EF{(b) t-SNE of the same states}{(b) t-SNE des mêmes états}},
  xlabel={\EF{neighbourhoods kept, distances not}{voisinages gardés, distances non}},
  label style={text=black!55}]
\addplot[lat/nuage gris] table[x=x, y=y, meta=marge] {\LatDonnees/latents-etats-tsne.csv};
\addplot[t3 rappel, quiver={u=\thisrow{tsnex}-\thisrow{tsnelx}, v=\thisrow{tsney}-\thisrow{tsnely}}, forget plot]
  table[x=tsnelx, y=tsnely] {\LatDonnees/latents-etats-reperes.csv};
\addplot[lat/noms] table[x=tsnelx, y=tsnely, meta=expert] {\LatDonnees/latents-etats-reperes.csv};
\end{axis}
% (c) les petits multiples
\node[t3 titre panneau] at (-0.02cm, -5.55cm) {\EF{(c) one panel per first expert, on the t-SNE map}%
  {(c) un panneau par expert premier, sur la carte de la t-SNE}};
\LatMultiple{1}{Un}{0}{0}\LatMultiple{2}{Deux}{1}{0}\LatMultiple{3}{Trois}{2}{0}\LatMultiple{4}{Quatre}{3}{0}%
\LatMultiple{5}{Cinq}{0}{1}\LatMultiple{6}{Six}{1}{1}\LatMultiple{7}{Sept}{2}{1}\LatMultiple{8}{Huit}{3}{1}%
}{}
% la clé, sous les panneaux
\node[anchor=north west, t3 sans cesure, font=\sffamily\scriptsize, align=left, inner sep=0pt, text=black!75,
  text width=12.2cm] at (0, -11.4cm) {%
  % troisième passe (N6 du juge latents, DN13 du juge de design) : chaque gris a sa pastille, et l'état fragile
  % de (c) son anneau
  \mbox{\LatPastille{latCarte}{1.6}\,\EF{(a), (b) a state}{(a), (b) un état}}\hspace{1.1em}%
  \mbox{\LatPastille{latFond}{1.6}\,\EF{(c) a state of another expert}{(c) un état d'un autre expert}}\hspace{1.1em}%
  \mbox{\LatPastille{latFragile}{1.6}\,\tikz[baseline=-0.55ex]\draw[latFragile, fill=latExpert!45, line width=0.6pt] (0,0) circle (1.9pt);\,%
    \EF{margin $z_{(1)}-z_{(2)}<\LatMargeDouzeDix$ (\LatMargeQuantile\%; ring in (c))}%
    {marge $z_{(1)}-z_{(2)}<\LatMargeDouzeDix$ (\LatMargeQuantile\,\% ; anneau dans (c))}}\hspace{1.1em}%
  \mbox{\LatPastille{latExpert}{1.6}\,\EF{(c) the panel's expert}{(c) l'expert du panneau}}\hspace{1.1em}%
  \mbox{\textbf{\LatGroupeUn}\enspace{\color{black!55}\LatGroupeUnPart\LatPc}\enspace\EF{expert, its share of the states}{expert, sa part des états}}};
\end{tikzpicture}}

\Titre{LatentsEtats}{The latent space of the \LatEtats{} states of layer~\LatCouche{} of gpt-oss-20b}%
{L'espace latent des \LatEtats{} états de la couche~\LatCouche{} de gpt-oss-20b}
\Legende{LatentsEtats}{\mes}%
{The \LatEtats{} states of layer~\LatCouche{} gather by first expert, and the states of small margin
stand more often on the edges of the groups (\LatBordFragile{} of their neighbours with another expert, against
\LatBordLarge).}%
{Les \LatEtats{} états de la couche~\LatCouche{} se regroupent par expert premier, et les états de
petite marge se tiennent plus souvent aux bords des groupes (\LatBordFragile{} de voisins d'un autre expert, contre
\LatBordLarge).}%
{Calibration states of gpt-oss-20b ($h$ after the post-attention norm, \LatDim{} dimensions, windows 10 to 21 of
wiki.test). First expert: the largest router logit; margin: $z_{(1)}-z_{(2)}$, the gap between the two largest biased
logits. In black (ringed in (c)), the \LatFragiles{} states of smallest margin (\LatMargeQuantile\%, below \LatMargeDouzeDix).
(a)~PCA, centred and unscaled: \LatAcpVarPlan\% of the variance, trustworthiness $T(\LatK)=\LatAcpT$, continuity
$C(\LatK)=\LatAcpC$, Spearman correlation of distances \LatAcpRho. (b)~t-SNE written from scratch, deterministic
(perplexity \LatTsnePerplexite): $T=\LatTsneT$, $C=\LatTsneC$, Spearman \LatTsneRho; kept by the threshold declared
before the run ($T$ and $C$ at least those of the PCA), which its PCA initialisation already meets: it exceeds it by
\LatTsneGainT{} and \LatTsneGainC. Its distances between groups and its areas do not read. Numbers: where the states
of the group are densest (density, not purity: in the PCA the group is often a minority there), moved aside a little
where two names overlapped. (c)~The eight most frequent first experts (\LatGroupesPart\% of the states;
\LatAutresExperts{} others). Share of the \LatK{} nearest neighbours with the same first expert: \LatAccordHaut\% in
the original space, \LatAcpAccord\% (PCA), \LatTsneAccord\% (t-SNE), \LatAccordHasard\% by chance (sum of the squared
shares). Edge of a state, the share of its \LatKBord{} neighbours with another first expert: \LatBordFragile{} for the
states of small margin against \LatBordLarge{} beyond the median (t-SNE plane), \LatBordOrigFragile{} against
\LatBordOrigLarge{} in the original space (judge of the first flow); the margin of the set of four experts served,
$z_{(4)}-z_{(5)}$, does not mark them (\LatBordQuatreCinq{} below the 10th centile, as \LatBordQuatreCinqLarge{} beyond the median) and is not drawn. Controls: the router re-read reproduces the
\LatWzeroEnsembles{} sets of the log (without its bias: \LatWzeroCassure); fidelity 1 for the full PCA,
$T=\LatTemoinPermuteT$ for a permuted plane; edges equal to the judge's (with \LatKBordCasse{} neighbours, they depart
from them).}%
{États de calibration de gpt-oss-20b ($h$ après la norme d'après l'attention, \LatDim{} dimensions, fenêtres 10 à 21
de wiki.test). Expert premier : le plus grand logit du routeur ; marge : $z_{(1)}-z_{(2)}$, l'écart des deux plus
grands logits biaisés. En noir (cerclés dans (c)), les \LatFragiles{} états de plus petite marge (\LatMargeQuantile\,\%, sous
\LatMargeDouzeDix). (a)~ACP centrée, non réduite : \LatAcpVarPlan\,\% de la variance, confiance $T(\LatK)=\LatAcpT$,
continuité $C(\LatK)=\LatAcpC$, Spearman des distances \LatAcpRho. (b)~t-SNE écrite de zéro, déterministe (perplexité
\LatTsnePerplexite) : $T=\LatTsneT$, $C=\LatTsneC$, Spearman \LatTsneRho{} ; gardée par le seuil déclaré avant la
passe ($T$ et $C$ au moins ceux de l'ACP), que son initialisation par l'ACP atteint déjà : elle le dépasse de
\LatTsneGainT{} et de \LatTsneGainC. Ses distances entre groupes et ses aires ne se lisent pas. Numéros : là où les
états du groupe sont le plus serrés (la densité, pas la pureté : dans l'ACP, le groupe y est souvent minoritaire), un
peu écartés là où deux noms se recouvraient. (c)~Les huit experts premiers les plus fréquents (\LatGroupesPart\,\% des états ;
\LatAutresExperts{} autres). Part des \LatK{} plus proches voisins de même expert premier : \LatAccordHaut\,\% dans
l'espace d'origine, \LatAcpAccord\,\% (ACP), \LatTsneAccord\,\% (t-SNE), \LatAccordHasard\,\% au hasard (somme des
carrés des parts). Bord d'un état, la part de ses \LatKBord{} voisins d'un autre expert premier : \LatBordFragile{}
pour les états de petite marge contre \LatBordLarge{} au-delà de la médiane (plan de la t-SNE), \LatBordOrigFragile{}
contre \LatBordOrigLarge{} dans l'espace d'origine (juge du premier flux) ; la marge de l'ensemble des quatre experts
servis, $z_{(4)}-z_{(5)}$, ne les marque pas (\LatBordQuatreCinq{} sous le 10\textsuperscript{e} centile comme \LatBordQuatreCinqLarge{} au-delà de la médiane) et n'est pas dessinée. Témoins : le routeur relu
redonne les \LatWzeroEnsembles{} ensembles du journal (sans son biais : \LatWzeroCassure) ; fidélité 1 pour l'ACP
complète, $T=\LatTemoinPermuteT$ pour un plan permuté ; bords égaux à ceux du juge (à \LatKBordCasse{} voisins, ils
s'en écartent).}
\Source{LatentsEtats}{\fichier{mesures/genome/tome3/latents/etats-c23.csv}, \fichier{etats-c23.json}
(\EF{seed}{graine} \LatGraine), \EF{by}{par} \fichier{outils/genome/tome3/latents_etats.py}\Pv{}
\EF{captures}{captures} \fichier{h-L23.f32}, \fichier{routage-L23.bin} (\EF{scout}{éclaireur})\Pv{}
\EF{edges in the original space}{bords dans l'espace d'origine} \fichier{outils/genome/tome3/juge-latents/juge-etats-marges.json}\Pv{}
\fichier{memoire/outils/tome3_latents.py} (\EF{controls}{témoins} \fichier{latents-temoins.json}).}

% =====================================================================================================
% III-9.1 (d). Une rangée de 32 poids : exact, frontière, synonyme ; les codons changés.
\newcommand{\LatRangeeLigne}[3]{% #1 ordonnée, #2 colonne du CSV, #3 fichier
\addplot[only marks, mark=none, point meta=explicit symbolic, nodes near coords, nodes near coords align=center,
  every node near coord/.append style={font=\scriptsize\sffamily, anchor=center, inner sep=0pt}]
  table[x=j, y expr=#1, meta=#2] {\LatDonnees/latents-#3.csv};}
\newcommand{\GrapheLatentsRangee}{%
\begin{tikzpicture}
\begin{axis}[t3 graphe, t3 points, width=0.93\linewidth, height=6.4cm, xmin=-0.7, xmax=31.7, ymin=-4.6, ymax=4.6,
  clip=false, xtick={0,...,31}, xticklabels={}, ytick={-4,-3,-2,-1,0,1,2,3,4}, minor ytick={-1.5,-0.5,0.5,1.5},
  grid=none, ymajorgrids, yminorgrids, major grid style={black!9}, minor grid style={black!5},
  tick style={black!30}, xtick style={draw=none}, tick label style={font=\scriptsize}, label style={font=\small},
  ylabel={\EF{weight / block scale $s_b$}{poids / échelle du bloc $s_b$}},
  title={\EF{(c) one row of 32 weights, before and after}{(c) une rangée de 32 poids, avant et après}},
  title style={yshift=0.62cm},
  legend style={at={(0.5,1.02)}, anchor=south, legend columns=4, draw=none, fill=none,
    /tikz/every even column/.append style={column sep=0.45cm}}]
% les colonnes dont le codon change
\addplot[ybar, bar width=0.86, bar shift=0pt, draw=none, fill=CouleurSynonyme!10, forget plot]
  table[x=j, y expr=4.6] {\LatDonnees/latents-rangee-change.csv};
\addplot[ybar, bar width=0.86, bar shift=0pt, draw=none, fill=CouleurSynonyme!10, forget plot]
  table[x=j, y expr=-4.6] {\LatDonnees/latents-rangee-change.csv};
% les huit niveaux de la base
\foreach \latniv in {\LatNivUn,\LatNivDeux,\LatNivTrois} {
  \edef\latcmd{\noexpand\draw[CouleurFrontiere, line width=0.5pt, dash pattern=on 2.2pt off 1.6pt]
    (axis cs:-0.7,\latniv) -- (axis cs:31.7,\latniv) (axis cs:-0.7,-\latniv) -- (axis cs:31.7,-\latniv);}\latcmd}
% de la base frontière à la base synonyme ; troisième passe (N5 du juge latents) : la flèche est posée à droite
% des marques, sans raccourci au départ, pour que la plus courte (±c_0, en j = 11) se voie aussi
\addplot[quiver={u=0, v={\thisrow{synonyme}-\thisrow{frontiere}}}, draw=CouleurSynonyme, line width=0.9pt,
  -{Stealth[length=1.4mm]}, shorten >=0.6pt, forget plot]
  table[x expr=\thisrow{j}+0.36, y=frontiere] {\LatDonnees/latents-rangee-change.csv};
\addplot[only marks, t3 frontiere point] table[x=j, y=frontiere] {\LatDonnees/latents-rangee.csv};
\addlegendentry{\EF{frontier base}{base frontière}}
\addplot[only marks, t3 synonyme point] table[x=j, y=synonyme] {\LatDonnees/latents-rangee.csv};
\addlegendentry{\EF{synonym base}{base synonyme}}
\addplot[only marks, mark=triangle*, mark size=1.9pt, mark options={solid, fill=CouleurExact, draw=white,
  line width=0.3pt}] table[x=j, y=exact] {\LatDonnees/latents-rangee.csv};
\addlegendentry{\EF{exact weight (E2M1 grid)}{poids exact (grille E2M1)}}
\addlegendimage{area legend, draw=none, fill=CouleurSynonyme!10}
\addlegendentry{\EF{codon changed (\LatRangChanges{} of 32)}{codon changé (\LatRangChanges{} sur 32)}}
% étiquettes des niveaux, à droite
\node[t3 etiquette, anchor=west, text=black!60] at (axis cs:31.8,0) {$\pm c_0$};
\node[t3 etiquette, anchor=west, text=black!60] at (axis cs:31.8,\LatNivDeux) {$+c_1$};
\node[t3 etiquette, anchor=west, text=black!60] at (axis cs:31.8,\LatNivTrois) {$+c_2$};
\node[t3 etiquette, anchor=west, text=black!60] at (axis cs:31.8,-\LatNivDeux) {$-c_1$};
\node[t3 etiquette, anchor=west, text=black!60] at (axis cs:31.8,-\LatNivTrois) {$-c_2$};
% la bande des codons
\LatRangeeLigne{-5.3}{j}{rangee}
\LatRangeeLigne{-6.1}{code}{rangee}
\LatRangeeLigne{-6.9}{sf}{rangee}
\LatRangeeLigne{-7.7}{ss}{rangee-garde}
\begin{scope}[every node near coord/.append style={text=CouleurSynonyme, font=\scriptsize\sffamily\bfseries}]
\LatRangeeLigne{-7.7}{ss}{rangee-change}
\end{scope}
\node[anchor=east, t3 sans cesure, font=\scriptsize\sffamily, text=black!55] at (axis cs:-0.8,-5.3) {$j$};
\node[anchor=east, t3 sans cesure, font=\scriptsize\sffamily, text=black!70] at (axis cs:-0.8,-6.1) {\EF{MXFP4 code}{code MXFP4}};
\node[anchor=east, t3 sans cesure, font=\scriptsize\sffamily, text=black!70] at (axis cs:-0.8,-6.9) {\EF{frontier $\pm k$}{frontière $\pm k$}};
\node[anchor=east, t3 sans cesure, font=\scriptsize\sffamily, text=CouleurSynonyme] at (axis cs:-0.8,-7.7) {\EF{synonym $\pm k$}{synonyme $\pm k$}};
\draw[black!15] (axis cs:-0.7,-5.7) -- (axis cs:31.7,-5.7);
\end{axis}
\end{tikzpicture}}
% les niveaux, lus du CSV (aucun nombre recopié)
\pgfplotstableread{\LatDonnees/latents-rangee-niveaux.csv}\LatNiveaux
\pgfplotstablegetelem{0}{y}\of{\LatNiveaux}\edef\LatNivUn{\pgfplotsretval}
\pgfplotstablegetelem{2}{y}\of{\LatNiveaux}\edef\LatNivDeux{\pgfplotsretval}
\pgfplotstablegetelem{4}{y}\of{\LatNiveaux}\edef\LatNivTrois{\pgfplotsretval}

\Titre{LatentsRangee}{A row of 32 weights before and after the synonymous mutation}%
{Une rangée de 32 poids avant et après la mutation synonyme}
\Legende{LatentsRangee}{\mes}%
{The synonymous mutation changes \LatRangChanges{} of the 32 symbols of this row and moves \LatRangEloignes{} of
these weights away from their exact value (the other one, a zero, changes sign at the same distance): it is
synonymous for the output, not for the weights.}%
{La mutation synonyme change \LatRangChanges{} des 32 symboles de cette rangée et en éloigne \LatRangEloignes{}
de leur valeur exacte (l'autre, un zéro, change de signe à la même distance) : elle est synonyme pour la sortie,
pas pour les poids.}%
{Panel (c) of \cref{fig:t3-syno-base-propre}: a row of 32 weights of the same matrix as panels (a), (b) and (b$'$) (gpt-oss-20b, layer~8,
expert~0, gate): row \LatRangLigne, block \LatRangBloc{} (columns \LatRangColA{} to \LatRangColB), in units of the
block scale $s_b=2^{\LatRangExpoRel}$. Black triangles: the exact weight, on the E2M1 grid; grey diamonds: the
frontier base (its symbol is a function of the code); blue circles: the synonym base ($\lambda=1$, $\delta^\ast=1$),
as served by the lot~B \texttt{.sym} files; an arrow, beside the marks, goes from the frontier value to the synonym value where the
symbol changes. Dashed lines: the levels $\pm c_k$ of the base ($c_0=\LatCentreUn$, $c_1=\LatCentreDeux$,
$c_2=\LatCentreTrois$; $c_3=\LatCentreQuatre$ is out of the frame), the same for both bases. In the band, $j$ is the
column in the block, then the MXFP4 code (hexadecimal), then each base's symbol: ``$+k$'' and ``$-k$'' stand for the
level $\pm c_k$ ($k=0$ to $3$), not for a signed zero. The block was drawn (seed) among the \LatRangCandidats{} blocks
whose number of changed symbols equals the median over the matrix (\LatRangMediane{} of 32; mean \LatRangMoyenne,
\LatRangPartMatrice\% of the symbols of the matrix): typical by construction. In this block the synonym raises the
weight error (\LatRangErrS{} against \LatRangErrF, in $s_b^2$); over the whole matrix the error energy rises by
\LatSpecEnergieHausse\% (\LatSpecEnergieF{} to \LatSpecEnergieS) while the output error falls
\LatSpecCalibRapport-fold on calibration (\LatSpecCalibF{} to \LatSpecCalibS) and \LatSpecEvalRapport-fold seen by
the evaluation. Not shown: the output, which only (a), (b) and (b$'$) show; one row cannot show where the error goes (to the
mute directions of $H$, panel (b)). Control: the symbols read reproduce $\lVert W-\hat W\rVert_F^2=\sum_k e_k$ of the
lot~B file to \LatRangTemoin{} relative (another expert's files: a relative gap of \LatRangCassure).}%
{Panneau (c) de la \cref{fig:t3-syno-base-propre} : une rangée de 32 poids de la même matrice que les panneaux (a), (b) et (b$'$) (gpt-oss-20b,
couche~8, expert~0, gate) : ligne \LatRangLigne, bloc \LatRangBloc{} (colonnes \LatRangColA{} à \LatRangColB), en
unités de l'échelle du bloc $s_b=2^{\LatRangExpoRel}$. Triangles noirs : le poids exact, sur la grille E2M1 ; losanges
gris : la base frontière (son symbole est une fonction du code) ; ronds bleus : la base synonyme ($\lambda=1$,
$\delta^\ast=1$), telle que la servent les fichiers \texttt{.sym} du lot~B ; une flèche, à côté des marques, va de la valeur frontière à la
valeur synonyme là où le symbole change. Tirets : les niveaux $\pm c_k$ de la base ($c_0=\LatCentreUn$,
$c_1=\LatCentreDeux$, $c_2=\LatCentreTrois$ ; $c_3=\LatCentreQuatre$ sort du cadre), les mêmes pour les deux bases.
Dans la bande, $j$ est la colonne dans le bloc, puis le code MXFP4 (hexadécimal), puis le symbole de chaque base :
« $+k$ » et « $-k$ » notent le niveau $\pm c_k$ ($k=0$ à $3$), pas un zéro signé. Le bloc est tiré (graine) parmi les
\LatRangCandidats{} blocs dont le nombre de symboles changés égale la médiane sur la matrice (\LatRangMediane{} sur 32 ;
moyenne \LatRangMoyenne, \LatRangPartMatrice\,\% des symboles de la matrice) : typique par construction. Dans ce bloc,
le synonyme augmente l'erreur sur les poids (\LatRangErrS{} contre \LatRangErrF, en $s_b^2$) ; sur la matrice entière,
l'énergie de l'erreur monte de \LatSpecEnergieHausse\,\% (\LatSpecEnergieF{} à \LatSpecEnergieS) quand l'erreur de
sortie baisse \LatSpecCalibRapport{} fois en calibration (\LatSpecCalibF{} à \LatSpecCalibS) et
\LatSpecEvalRapport{} fois vue par l'évaluation. Non montré : la sortie, que seuls (a), (b) et (b$'$) montrent ; une rangée ne
peut pas montrer où va l'erreur (vers les directions muettes de $H$, panneau (b)). Témoin : les symboles lus redonnent
$\lVert W-\hat W\rVert_F^2=\sum_k e_k$ du fichier du lot~B à \LatRangTemoin{} relatif près (les fichiers d'un autre
expert : un écart relatif de \LatRangCassure).}
\Source{LatentsRangee}{\fichier{mesures/genome/tome3/latents/rangee-L08-E00-gate.csv}, \fichier{rangee-L08-E00-gate.json}
(\EF{seed}{graine} \LatRangGraine), \EF{by}{par} \fichier{outils/genome/tome3/latents_rangee.py}\Pv{}
\fichier{mesures/genome/construction/synonyme/figures/spectre-L08-E00-gate.csv}\Pv{}
\fichier{memoire/outils/tome3_latents.py}.}

% =====================================================================================================
% III-9.7. L'erreur de sortie d'un expert dans les directions de ses entrées : la fonction d'appui
% θ ↦ ‖ΔW v(θ)‖ des deux bases dans les plans (u1, u2) et (u1, u3), la rose des directions des entrées, et
% le rapport synonyme / frontière le long de u1, u2, u3.
% Un panneau polaire. #1 : plan (12, 13) ; #2 : indice du second axe (2, 3) ; #3 : options.
% Du centre vers le bord : la courbe synonyme, la rose des entrées (en anneau, sans unité radiale), la courbe
% frontière. L'échelle radiale (unités de la sortie) se lit sur les graduations du demi-axe bas.
\newcommand{\LatPanneauAppui}[3]{%
\begin{axis}[t3 graphe, scale only axis, width=4.6cm, height=4.6cm, axis equal image, axis lines=none,
  xmin=-\LatNuageRayon, xmax=\LatNuageRayon, ymin=-\LatNuageRayon, ymax=\LatNuageRayon, clip=false, #3]
% les deux axes propres, et la graduation radiale sur le demi-axe bas
\draw[black!30, line width=0.4pt] (axis cs:-\LatNuageRayon,0) -- (axis cs:\LatNuageRayon,0)
  (axis cs:0,-\LatNuageRayon) -- (axis cs:0,\LatNuageRayon);
\node[t3 etiquette, anchor=west, text=black!65, inner sep=1.5pt] at (axis cs:\LatNuageRayon,0) {$u_1$};
\node[t3 etiquette, anchor=south, text=black!65, inner sep=1.5pt] at (axis cs:0,\LatNuageRayon) {$u_#2$};
\foreach \latr in {0.05,0.10} {
  \edef\latcmd{\noexpand\draw[black!45, line width=0.4pt] (axis cs:-0.004,-\latr) -- (axis cs:0.004,-\latr);
    \noexpand\node[t3 etiquette, anchor=west, text=black!55, inner sep=0.8pt] at (axis cs:0.005,-\latr)
      {\noexpand\pgfmathprintnumber[fixed, fixed zerofill, precision=2]{\latr}};}\latcmd}
% la rose des directions des entrées
\addplot[quiver={u=\thisrow{dx}, v=\thisrow{dy}}, draw=latEntree, line width=1.5pt, line cap=butt, forget plot]
  table[x=x0, y=y0] {\LatDonnees/latents-rose-#1.csv};
% les deux fonctions d'appui
\addplot[t3 frontiere, forget plot] table[x=x, y=y] {\LatDonnees/latents-appui-frontiere-#1.csv};
\addplot[t3 synonyme, fill=CouleurSynonyme, fill opacity=0.18, forget plot]
  table[x=x, y=y] {\LatDonnees/latents-appui-synonyme-#1.csv};
\fill[Encre] (axis cs:0,0) circle (0.8pt);
\end{axis}}
\newcommand{\GrapheLatentsNuage}{%
\begin{tikzpicture}
\LatPanneauAppui{12}{2}{name=latappuia, at={(0,0)}, anchor=south west,
  title={\EF{(a) plane of axes 1 and 2 of $H$}{(a) plan des axes 1 et 2 de $H$}}}
\LatPanneauAppui{13}{3}{name=latappuib, at={(5.15cm,0)}, anchor=south west,
  title={\EF{(b) plane of axes 1 and 3}{(b) plan des axes 1 et 3}}}
% (c) le rapport, axe par axe
% Troisième passe (GN1 du juge latents, DN15 du juge de design) : une quatrième barre, à droite d'un filet, donne le
% même rapport sur les entrées réelles entières (médiane, et leur étendue en moustache) : c'est elle que la tête cite.
\begin{axis}[t3 graphe, ybar, t3 barres, name=latbarres, at={(11.15cm,0)}, anchor=south west, scale only axis,
  width=3.25cm, height=4.6cm, bar width=0.40cm, xmin=0.45, xmax=4.85, ymin=0, ymax=1.18,
  xtick={1,2,3,4.25}, xticklabels={$u_1$,$u_2$,$u_3$,$x$}, xtick style={draw=none},
  ytick={0,0.5,1}, t3 decimales y=1, grid=none, ymajorgrids, major grid style={black!8},
  tick label style={font=\scriptsize}, label style={font=\scriptsize},
  ylabel={\EF{synonym / frontier}{synonyme / frontière}}, ylabel shift=-2pt,
  title={\EF{(c) along each axis}{(c) axe par axe}},
  bar shift=0pt, clip=false]
\addplot[fill=CouleurSynonyme, draw=CouleurSynonyme!70!black, line width=0.4pt,
  nodes near coords={\pgfmathprintnumber[fixed, fixed zerofill, precision=3]{\pgfplotspointmeta}},
  nodes near coords style={t3 nombre=CouleurSynonyme, font=\sffamily\scriptsize, anchor=south, inner sep=1pt}]
  table[x=k, y=rapport] {\LatDonnees/latents-barres.csv};
\addplot[fill=CouleurSynonyme!35, draw=CouleurSynonyme!70!black, line width=0.4pt,
  error bars/.cd, y dir=both, y explicit, error mark=-, error bar style={Encre!75, line width=0.6pt}]
  table[x=k, y=rapport, y error plus=haut, y error minus=bas] {\LatDonnees/latents-barres-entrees.csv};
\draw[black!30, line width=0.4pt] (axis cs:3.6,0) -- (axis cs:3.6,1.18);
\node[t3 nombre=CouleurSynonyme, anchor=south, inner sep=1pt] at (axis cs:4.25,\LatNuageEntreesMaxBrut) {\LatNuageEntreesMed};
\draw[t3 frontiere] (axis cs:0.45,1) -- (axis cs:4.85,1);
\node[t3 nombre=CouleurFrontiereMarque, anchor=south east, inner sep=1pt] at (axis cs:3.55,1.01) {\EF{frontier}{frontière}};
% λ_k sous chaque axe, en gris : où les entrées pèsent
\node[t3 sans cesure, font=\sffamily\scriptsize, text=black!55, anchor=east, inner sep=0pt] at (axis cs:0.4,-0.17) {$\lambda_k$};
\node[t3 sans cesure, font=\sffamily\scriptsize, text=black!55, anchor=base, inner sep=0pt] at (axis cs:1,-0.2) {\LatNuageLambdaUn};
\node[t3 sans cesure, font=\sffamily\scriptsize, text=black!55, anchor=base, inner sep=0pt] at (axis cs:2,-0.2) {\LatNuageLambdaDeux};
\node[t3 sans cesure, font=\sffamily\scriptsize, text=black!55, anchor=base, inner sep=0pt] at (axis cs:3,-0.2) {\LatNuageLambdaTrois};
\end{axis}
% la clé, sous les panneaux
\node[anchor=north west, t3 sans cesure, font=\sffamily\scriptsize, align=left, inner sep=0pt, text=black!75,
  text width=13.8cm] at ([yshift=-0.55cm]latappuia.south west) {%
  \mbox{\tikz[baseline=-0.55ex]\draw[t3 frontiere] (0,0) -- (0.55,0);\,%
    \EF{frontier: $\lVert\Delta W v\rVert$ for a unit input $v$}{frontière : $\lVert\Delta W v\rVert$ pour une entrée unité $v$}}\hspace{1.2em}%
  \mbox{\tikz[baseline=-0.55ex]{\fill[CouleurSynonyme, fill opacity=0.18] (0,-0.09) rectangle (0.55,0.09);
    \draw[t3 synonyme] (0,0) -- (0.55,0);}\,\EF{synonym, the same}{synonyme, de même}}\hspace{1.2em}%
  % troisième passe (DN15) : la rose a une icône de barre, pas de trait (elle se confondait avec la frontière)
  \mbox{\tikz[baseline=-0.55ex]{\fill[latEntree] (0,-0.07) rectangle (0.12,0.07);
    \fill[latEntree] (0.16,-0.07) rectangle (0.24,0.04);\fill[latEntree] (0.28,-0.07) rectangle (0.34,0.0);}\,%
    \EF{projected directions of the \LatNuageN{} inputs, per \LatNuageRosePas$^\circ$ slice;}{directions projetées des
      \LatNuageN{} entrées, par tranche de \LatNuageRosePas$^\circ$ ;}}\space
  \mbox{\EF{median angle of an input to the plane $(u_1, u_2)$: \LatNuageAngleMed$^\circ$}{angle médian d'une entrée
      au plan $(u_1, u_2)$ : \LatNuageAngleMed$^\circ$}}};
\end{tikzpicture}}

\Titre{LatentsNuage}{The output error of an expert in the directions of its inputs, frontier and synonym}%
{L'erreur de sortie d'un expert dans les directions de ses entrées, frontière et synonyme}
\Legende{LatentsNuage}{\mes\,\calc}%
{In the planes of the first axes of $H$, where the expert's inputs project, the synonym brings the output error
down to at most \LatNuageConeRapport{} of the frontier's; on the whole inputs, \LatNuageHorsTrois\% of whose
squared norm lies outside these three axes, to \LatNuageEntreesMed{} (median; \LatNuageEntreesMin{} to
\LatNuageEntreesMax).}%
{Dans les plans des premiers axes de $H$, où se projettent les entrées de l'expert, le synonyme ramène l'erreur de
sortie à au plus \LatNuageConeRapport{} de celle de la frontière ; sur les entrées entières, dont \LatNuageHorsTrois\,\%
de la norme au carré est hors de ces trois axes, à \LatNuageEntreesMed{} (médiane ; de \LatNuageEntreesMin{} à
\LatNuageEntreesMax).}%
{Expert~0 of layer~8 of gpt-oss-20b, gate matrix, and its \LatNuageN{} real calibration inputs $x$ (the positions the
router sends it, windows 10 to 21 of wiki.test), on the first eigen-axes $u_k$ of $H=\sum g^2xx^\top/\sum g^2$ (not
damped, the $H$ of \cref{fig:t3-syno-base-propre}; $\lambda_1=\LatNuageLambdaUn$, $\lambda_2=\LatNuageLambdaDeux$,
$\lambda_3=\LatNuageLambdaTrois$, \LatNuagePartTrace\% of the trace). (a), (b)~In each unit direction $v$ of the plane,
each closed curve reaches $\lVert\Delta W v\rVert$, the output error for a unit input of direction $v$, with $\Delta
W=W-\hat W$ for that base: the support function of the ellipse $M^{1/2}$, $M=B^\top\Delta W^\top\Delta W B$,
$B=[u_1\ u_k]$; the radial scale, in the units of the output, is marked on the lower half-axis. The frontier is almost a
circle: it spreads its error evenly. In the ring between the two curves, the rose counts the projected directions of
the inputs, their projections onto the plane (bar length proportional to the count; its radial place carries no unit):
\LatNuageConeQuantile\% of the projections lie within $\pm\LatNuageConeDouze^\circ$ of $u_1$ in (a) and
$\pm\LatNuageConeTreize^\circ$ in (b); for a unit input in that cone, the synonym's error is at most
\LatNuageConeRapportDouze{} (a) and \LatNuageConeRapportTreize{} (b) of the frontier's. $H$ is not centred: $u_1$ carries the mean of the inputs, all near
$u_1^\top x=\LatNuageMoyUn$ (spread \LatNuageEcartUn, against \LatNuageEcartDeux{} along $u_2$), and $\lambda_1$ owes
most to that mean ($\LatNuageMoyUn^2\approx\LatNuageMoyCarre$). (c)~The ratio synonym / frontier of $\lVert\Delta W
u_k\rVert$ along each axis: \LatNuageRapportUn{} along $u_1$, the axis where the inputs weigh most (their mean),
\LatNuageRapportDeux{} along $u_2$, \LatNuageRapportTrois{} along $u_3$ (frontier \LatNuageErrFUn,
\LatNuageErrFDeux, \LatNuageErrFTrois). Right of the rule, the same ratio for the \LatNuageN{} real inputs in their
\LatDim{} dimensions, $\lVert\Delta W_s x\rVert/\lVert\Delta W_f x\rVert$: median \LatNuageEntreesMed, from
\LatNuageEntreesMin{} to \LatNuageEntreesMax{} (whisker), \LatNuageEntreesSousTete\% of them below
\LatNuageConeRapport; an input makes a median angle of \LatNuageAngleMed$^\circ$ with the plane $(u_1, u_2)$. These
ratios come from the journal of the second-flow judge, recomputed from the GGUF, the \texttt{.sym} files and the
captures; a control ties them to the committed files: their $g^2$-weighted quadratic ratio equals
$\sqrt{\sum_k\lambda_k e_k^{\mathrm{syn}}/\sum_k\lambda_k e_k^{\mathrm{fr}}}$ of the lot~B file to $10^{-9}$ (with
$\lambda_1$ off by 1\%, it departs). Not shown: the \LatNuageAutresDirections{} other directions, where the synonym
puts the weight error it removes here (panel (b) of \cref{fig:t3-syno-base-propre}); weighted by the inputs, the output error falls there
too, from \LatNuageAutresF{} to \LatNuageAutresS, and these directions carry \LatNuageAutresPartGain\% of the output
gain; the held-out evaluation inputs, on which the gain shrinks (\LatSpecEvalRapport-fold against \LatSpecCalibRapport-fold). The curves are computed from the ellipses
committed by the lot: $M$ is recovered from their \LatNuageEllipsePoints{} points by least squares and recovers $e_k$ and the half-axes of
the lot's file to \LatNuageAppuiTemoin{} relative (with a scale factor off by \LatNuageCasseF\%, it departs from them). Control of the
lot: $\lambda_1$ to $\lambda_3$ and $e_1$ to $e_3$ equal those of the lot~B file to \LatNuageTemoin{} relative (from
the evaluation inputs, $\lambda_1$ moves by \LatNuageCassure\%).}%
{Expert~0 de la couche~8 de gpt-oss-20b, matrice gate, et ses \LatNuageN{} entrées réelles de calibration $x$ (les
positions que le routeur lui envoie, fenêtres 10 à 21 de wiki.test), sur les premiers axes propres $u_k$ de
$H=\sum g^2xx^\top/\sum g^2$ (non amortie, la $H$ de la \cref{fig:t3-syno-base-propre} ; $\lambda_1=\LatNuageLambdaUn$,
$\lambda_2=\LatNuageLambdaDeux$, $\lambda_3=\LatNuageLambdaTrois$, \LatNuagePartTrace\,\% de la trace). (a), (b)~Dans
chaque direction unité $v$ du plan, chaque courbe fermée atteint $\lVert\Delta W v\rVert$, l'erreur de sortie pour une
entrée unité de direction $v$, avec $\Delta W=W-\hat W$ pour cette base : la fonction d'appui de l'ellipse $M^{1/2}$,
$M=B^\top\Delta W^\top\Delta W B$, $B=[u_1\ u_k]$ ; l'échelle radiale, dans les unités de la sortie, est graduée sur le
demi-axe bas. La frontière est presque un cercle : elle répartit son erreur à plat. Dans l'anneau entre les deux
courbes, la rose compte les directions projetées des entrées, leurs projections sur le plan (longueur proportionnelle
au compte ; sa place radiale n'a pas d'unité) : \LatNuageConeQuantile\,\% des projections sont à moins de
$\pm\LatNuageConeDouze^\circ$ de $u_1$ dans (a), de $\pm\LatNuageConeTreize^\circ$ dans (b) ; pour une entrée unité de
ce cône, l'erreur du synonyme vaut au plus \LatNuageConeRapportDouze{} (a) et \LatNuageConeRapportTreize{} (b) de celle
de la frontière.
$H$ n'est pas centrée : $u_1$ porte la moyenne des entrées, toutes près de $u_1^\top x=\LatNuageMoyUn$ (écart-type
\LatNuageEcartUn, contre \LatNuageEcartDeux{} le long de $u_2$), et $\lambda_1$ doit l'essentiel à cette moyenne
($\LatNuageMoyUn^2\approx\LatNuageMoyCarre$). (c)~Le rapport synonyme / frontière de $\lVert\Delta W u_k\rVert$ le long
de chaque axe : \LatNuageRapportUn{} le long de $u_1$, l'axe où les entrées pèsent le plus (leur moyenne),
\LatNuageRapportDeux{} le long de $u_2$, \LatNuageRapportTrois{} le long de $u_3$ (frontière \LatNuageErrFUn,
\LatNuageErrFDeux, \LatNuageErrFTrois). À droite du filet, le même rapport pour les \LatNuageN{} entrées réelles dans
leurs \LatDim{} dimensions, $\lVert\Delta W_s x\rVert/\lVert\Delta W_f x\rVert$ : médiane \LatNuageEntreesMed, de
\LatNuageEntreesMin{} à \LatNuageEntreesMax{} (moustache), \LatNuageEntreesSousTete\,\% d'entre elles sous
\LatNuageConeRapport ; une entrée fait un angle médian de \LatNuageAngleMed$^\circ$ avec le plan $(u_1, u_2)$. Ces
rapports viennent du journal du juge du second flux, refaits depuis le GGUF, les fichiers \texttt{.sym} et les
captures ; un témoin les relie aux fichiers commités : leur rapport quadratique pondéré par $g^2$ égale
$\sqrt{\sum_k\lambda_k e_k^{\mathrm{syn}}/\sum_k\lambda_k e_k^{\mathrm{fr}}}$ du fichier du lot~B à $10^{-9}$ près
(avec $\lambda_1$ faux de 1\,\%, il s'en écarte). Non montrés : les \LatNuageAutresDirections{} autres directions, où le
synonyme met l'erreur sur les poids qu'il retire ici (panneau (b) de la \cref{fig:t3-syno-base-propre}) ; pondérée par les entrées, l'erreur de
sortie y baisse aussi, de \LatNuageAutresF{} à \LatNuageAutresS, et ces directions portent \LatNuageAutresPartGain\,\%
du gain de sortie ; les entrées d'évaluation tenues à l'écart, sur lesquelles le
gain se réduit (\LatSpecEvalRapport{} fois contre \LatSpecCalibRapport). Les courbes sont calculées depuis les
ellipses commitées par le lot : $M$ se retrouve de leurs \LatNuageEllipsePoints{} points par moindres carrés et redonne les $e_k$ et les
demi-axes du fichier du lot à \LatNuageAppuiTemoin{} relatif près (avec un facteur d'échelle faux de \LatNuageCasseF\,\%, elle s'en
écarte). Témoin du lot : $\lambda_1$ à $\lambda_3$ et $e_1$ à $e_3$ égalent ceux du fichier du lot~B à
\LatNuageTemoin{} relatif près (tirée des entrées d'évaluation, $\lambda_1$ bouge de \LatNuageCassure\,\%).}
\Source{LatentsNuage}{\fichier{mesures/genome/tome3/latents/nuage-L08-E00-gate-points.csv},
\fichier{nuage-L08-E00-gate-ellipses.csv}, \fichier{nuage-L08-E00-gate.json}, \EF{by}{par}
\fichier{outils/genome/tome3/latents_nuage.py}\Pv{} \EF{real inputs}{entrées réelles}
\fichier{outils/genome/tome3/juge2-latents/journaux/juge2-nuage-rangee.json}
(\texttt{entrees\_reelles})\Pv{} \fichier{memoire/outils/tome3_latents.py}
(\EF{controls}{témoins} \fichier{latents-temoins.json}).}
\fi
